\documentclass[12pt,a4paper]{article}

\usepackage[margin=2.5cm]{geometry}
\usepackage{amsmath, amssymb}
\usepackage{hyperref}
\usepackage{xcolor}
\usepackage{listings}
\usepackage{booktabs}
\usepackage{array}
\usepackage{longtable}
\usepackage{enumitem}
\usepackage{fancyhdr}
\usepackage{titlesec}
\usepackage{mdframed}
\usepackage[T1]{fontenc}
\usepackage[utf8]{inputenc}

\usepackage{tcolorbox}
\usepackage{multicol}
\usepackage{float}

\tcbuselibrary{breakable}

% Colors
\definecolor{angkorblue}{RGB}{0,70,127}
\definecolor{angkorgold}{RGB}{180,140,0}
\definecolor{codegreen}{RGB}{0,120,0}
\definecolor{codegray}{RGB}{100,100,100}
\definecolor{codeblue}{RGB}{0,0,180}
\definecolor{lightgray}{RGB}{248,248,248}
\definecolor{lessonbg}{RGB}{232,244,255}
\definecolor{warningbg}{RGB}{255,248,220}
\definecolor{successbg}{RGB}{232,255,232}
\definecolor{physicsbg}{RGB}{255,232,240}

% Code style
\lstdefinestyle{python}{
    backgroundcolor=\color{lightgray},
    commentstyle=\color{codegreen}\itshape,
    keywordstyle=\color{codeblue}\bfseries,
    stringstyle=\color{red!70!black},
    basicstyle=\ttfamily\footnotesize,
    breaklines=true,
    numbers=left,
    numberstyle=\tiny\color{codegray},
    frame=single,
    rulecolor=\color{angkorblue!40},
    language=Python,
    showstringspaces=false,
    tabsize=4,
    xleftmargin=15pt
}
\lstset{style=python}

% Box styles
\newtcolorbox{lessonbox}[1]{
    colback=lessonbg, colframe=angkorblue,
    fonttitle=\bfseries\color{white},
    title=#1, breakable
}
\newtcolorbox{physicsbox}[1]{
    colback=physicsbg, colframe=red!60!black,
    fonttitle=\bfseries\color{white},
    title=#1, breakable
}
\newtcolorbox{keybox}[1]{
    colback=successbg, colframe=green!50!black,
    fonttitle=\bfseries\color{white},
    title=#1, breakable
}
\newtcolorbox{warnbox}[1]{
    colback=warningbg, colframe=angkorgold,
    fonttitle=\bfseries\color{black},
    title=#1, breakable
}

% Headers
\pagestyle{fancy}
\fancyhf{}
\rhead{\textcolor{angkorblue}{ANGKOR -- Lesson Notes}}
\lhead{\textcolor{angkorblue}{\leftmark}}
\rfoot{\thepage}
\lfoot{\textcolor{codegray}{MUTH Boravy -- ITC Cambodia}}

\titleformat{\section}{\Large\bfseries\color{angkorblue}}{\thesection}{1em}{}[\titlerule]
\titleformat{\subsection}{\large\bfseries\color{angkorblue!80}}{\thesubsection}{1em}{}
\titleformat{\subsubsection}{\normalsize\bfseries\color{angkorblue!60}}{\thesubsubsection}{1em}{}

\hypersetup{colorlinks=true, linkcolor=angkorblue, urlcolor=angkorgold}

\begin{document}

% Title Page
\begin{titlepage}
\centering
\vspace*{1.5cm}
{\Huge\bfseries\color{angkorblue} ANGKOR}\\[0.4cm]
{\Large\color{angkorgold} Advanced Neutron Group-diffusion K-eigenvalue Of Reactors}\\[1cm]
\rule{\linewidth}{0.5mm}\\[0.3cm]
{\LARGE\bfseries Complete Lesson Notes \& Teaching Guide}\\[0.2cm]
{\large From Python Basics to International Benchmark Validation}\\
\rule{\linewidth}{0.5mm}\\[1.5cm]

\begin{minipage}{0.7\textwidth}\centering
{\large\textbf{Student:} MUTH Boravy}\\
Nuclear Engineering Student\\
Institute of Technology of Cambodia (ITC)\\
Phnom Penh, Cambodia\\[0.3cm]
\texttt{\href{https://github.com/muthboravy007/ANGKOR}{github.com/muthboravy007/ANGKOR}}
\end{minipage}

\vfill
{\large\itshape ``Named after Angkor Wat ---\\
because Cambodia's scientific legacy continues.''}\\[1cm]
{\large March 2026}
\end{titlepage}

\tableofcontents
\newpage

%==============================================================================
\section{Student Profile and Learning Preferences}
%==============================================================================

\begin{lessonbox}{How to Teach Boravy Effectively}
\begin{itemize}
  \item \textbf{Method:} Fill-in-the-blank exercises -- never give full code directly
  \item \textbf{Strength:} Strong physics intuition; uses physical reasoning to catch bugs
  \item \textbf{Style:} Visual learners -- tables, diagrams, side-by-side comparisons work best
  \item \textbf{Environment:} Windows, VS Code, Anaconda/venv, runs from ANGKOR root in terminal
  \item \textbf{Prior knowledge:} MCNP experience -- nuclear concepts transfer well
  \item \textbf{Docker:} Has OpenMC in Docker (openmc\_with\_jupyter image)
  \item \textbf{Pace:} One concept at a time; validate (run code) before advancing
\end{itemize}
\end{lessonbox}

\begin{warnbox}{Teaching Rules Established}
\begin{enumerate}
  \item Always ask the student to answer a physics question BEFORE coding it
  \item Fix bugs one at a time -- never all at once
  \item After every working result, ask ``why does this make sense physically?''
  \item If student copies code without understanding, return to blank-filling
\end{enumerate}
\end{warnbox}

%==============================================================================
\section{Lesson 1: Python Classes -- The Foundation}
%==============================================================================

\subsection{Concept: What is a Class?}

\begin{physicsbox}{Core Concept}
A class is a \textbf{blueprint}. From one blueprint, you can build many objects.
Each object has its own data (attributes) and behaviors (methods).
\end{physicsbox}

\subsection{Exercise 1: The Dog Class}

\begin{lessonbox}{Teaching Goal}
Introduce \texttt{\_\_init\_\_}, \texttt{self}, methods, and return values
using a non-physics example before applying to reactor geometry.
\end{lessonbox}

\begin{lstlisting}[caption={Dog Class -- First Exercise}]
class Dog:
    def __init__(self, name, breed, age):
        self.name  = name    # store name on the object
        self.breed = breed   # store breed
        self.age   = age     # store age

    def bark(self):
        print(f"{self.name} says: Woof!")

    def is_puppy(self):
        return self.age < 2  # True if under 2 years old

# --- Test ---
rex = Dog(name="Rex", breed="Labrador", age=1)
rex.bark()           # Rex says: Woof!
print(rex.is_puppy()) # True
\end{lstlisting}

\begin{keybox}{Key Lessons Learned}
\begin{itemize}
  \item \texttt{self.name = name} stores the input into the object
  \item Methods with no return give actions; methods with \texttt{return} give values
  \item Alignment formatting (\texttt{self.name\ \ \ = name}) is professional style
  \item Student scored 5/5 blanks correct on first attempt
\end{itemize}
\end{keybox}

\subsection{Exercise 2: The Region Class}

\begin{physicsbox}{Connection to Reactor Physics}
A reactor region is a rectangle with a material. The key question it must
answer: ``Is point (x,y) inside me?'' This maps directly to the Dog class pattern.
\end{physicsbox}

\begin{lstlisting}[caption={Region Class -- Applied to Reactor Geometry}]
class Region:
    """A rectangular region in 2D space with an assigned material."""

    def __init__(self, name, x_min, x_max, y_min, y_max, material):
        self.name     = name
        self.x_min    = x_min
        self.x_max    = x_max
        self.y_min    = y_min
        self.y_max    = y_max
        self.material = material
        self._validate()  # check immediately on creation

    def _validate(self):
        """Private method -- called automatically, not by user."""
        if self.x_min >= self.x_max:
            raise ValueError(
                f"Region '{self.name}': x_min must be < x_max")
        if self.y_min >= self.y_max:
            raise ValueError(
                f"Region '{self.name}': y_min must be < y_max")

    def contains(self, x, y):
        """Return True if point (x,y) is inside this region."""
        x_inside = self.x_min <= x <= self.x_max
        y_inside = self.y_min <= y <= self.y_max
        return x_inside and y_inside

    def width(self):
        return self.x_max - self.x_min   # DRY: defined once

    def height(self):
        return self.y_max - self.y_min

    def area(self):
        return self.width() * self.height()  # reuses width/height
\end{lstlisting}

\begin{keybox}{Design Decisions Made}
\begin{enumerate}
  \item \textbf{Cell-center mesh (Option 2):} Cell centers never land on boundaries,
        so \texttt{contains()} never gets an ambiguous boundary point.
        Used by CASMO, SIMULATE, PARCS.
  \item \textbf{Private \_validate():} Called in \texttt{\_\_init\_\_} automatically.
        User never needs to call it. Underscore = ``internal use only.''
  \item \textbf{DRY principle:} \texttt{area()} calls \texttt{width()} and
        \texttt{height()} -- if the formula changes, fix in one place only.
\end{enumerate}
\end{keybox}

%==============================================================================
\section{Lesson 2: 2D Geometry Engine}
%==============================================================================

\subsection{Concept: The Material Map}

\begin{physicsbox}{Core Physics Concept}
The geometry engine answers one fundamental question at every point in space:
\textbf{``What material is here?''} The answer is stored as a 2D array called
the \textbf{material map}.
\end{physicsbox}

\begin{lstlisting}[caption={Material Map -- 2D Array of Material Names}]
# Example material map for a 10x10 mesh:
# ['water' 'water' 'uo2_fuel' 'uo2_fuel' ... 'water']
# ['water' 'water' 'uo2_fuel' 'uo2_fuel' ... 'water']
# ...

# Built using cell centers:
dx = domain_x / nx          # cell width
dy = domain_y / ny          # cell height

x_centers = [(i + 0.5)*dx for i in range(nx)]  # never on boundary
y_centers = [(j + 0.5)*dy for j in range(ny)]

# Fill material map:
material_map = np.empty((ny, nx), dtype=object)  # rows=y, cols=x
for j in range(ny):
    for i in range(nx):
        x = x_centers[i]
        y = y_centers[j]
        for region in regions:
            if region.contains(x, y):
                material_map[j][i] = region.material
                break
\end{lstlisting}

\begin{warnbox}{Common Bug: Rows = Y, Columns = X}
NumPy arrays use \texttt{array[row][col]} = \texttt{array[j][i]}.
Rows go \textbf{down} (y-direction). Columns go \textbf{right} (x-direction).
Shape must be \texttt{(ny, nx)}, not \texttt{(nx, ny)}.
\end{warnbox}

\subsection{Cell Index Formula}

\begin{physicsbox}{Flattening 2D to 1D}
For a mesh with \texttt{nx} columns, cell $(i,j)$ maps to 1D index:
\begin{equation}
n = j \cdot n_x + i
\end{equation}
Moving \textbf{east} (i+1): $n$ increases by 1.
Moving \textbf{north} (j+1): $n$ increases by $n_x$ (one full row).
\end{physicsbox}

\begin{lstlisting}[caption={Cell Index -- Student Exercise that Built Understanding}]
# For a 4x3 mesh (nx=4, ny=3):
#      i=0   i=1   i=2   i=3
# j=2 [ 8  ][ 9  ][ 10 ][ 11 ]
# j=1 [ 4  ][ 5  ][ 6  ][ 7  ]
# j=0 [ 0  ][ 1  ][ 2  ][ 3  ]

# Cell (i=2, j=1):  n = 1*4 + 2 = 6
# East  (i=3, j=1): n = 7  = 6+1    (same row, +1)
# North (i=2, j=2): n = 10 = 6+4    (next row, +nx)
# West  (i=1, j=1): n = 5  = 6-1
# South (i=2, j=0): n = 2  = 6-4
\end{lstlisting}

\subsection{Geometry Visualizer}

\begin{lstlisting}[caption={Drawing Reactor Geometry with matplotlib}]
import matplotlib.patches as patches
import matplotlib.colors as mcolors

class GeometryVisualizer:
    MATERIAL_COLORS = {
        "water":       "steelblue",
        "uo2_fuel":    "orange",
        "control_rod": "dimgray",
    }

    def plot_geometry(self, show_mesh=False):
        fig, ax = plt.subplots(figsize=(10, 8))
        for region in self.engine.regions:
            color = self.MATERIAL_COLORS.get(region.material, "red")
            rect  = patches.Rectangle(
                (region.x_min, region.y_min),
                region.width(), region.height(),
                linewidth=2, edgecolor="black",
                facecolor=color, alpha=0.7
            )
            ax.add_patch(rect)
        ax.set_aspect("equal")
        plt.show()
\end{lstlisting}

%==============================================================================
\section{Lesson 3: YAML Input System}
%==============================================================================

\subsection{Why YAML?}

\begin{lessonbox}{Design Rationale}
Without YAML, users must edit Python code to change geometry -- that is a script,
not a tool. With YAML, users write a text file like MCNP input cards.
ANGKOR reads it and builds all objects automatically.
\end{lessonbox}

\begin{lstlisting}[language=Python, caption={How YAML Maps to Python Objects}]
# YAML file:
# materials:
#   uo2_fuel:
#     D1: 1.40
#     D2: 0.37

# Python access (after yaml.safe_load):
data["materials"]["uo2_fuel"]["D1"]  # -> 1.40

# Regions are a LIST (dash = list item):
# regions:
#   - name: fuel
#     x_min: 20.0
# Access:
for r in data["regions"]:
    name  = r["name"]    # "fuel"
    x_min = r["x_min"]  # 20.0
\end{lstlisting}

\begin{lstlisting}[caption={Auto-Detection of XS Format}]
def _parse_materials(self):
    """Reads YAML and auto-detects 2-group vs multi-group."""
    for mat_name, props in mats.items():
        if "groups" in props:
            # MULTI-GROUP FORMAT
            self.materials[mat_name] = {
                "groups":      props["groups"],
                "D":           props["D"],      # list of G values
                "sigma_a":     props["sigma_a"],
                "nu_sigma_f":  props["nu_sigma_f"],
                "sigma_s":     props["sigma_s"], # G x G matrix
                "chi":         props["chi"],
            }
        else:
            # 2-GROUP FORMAT
            self.materials[mat_name] = {
                "D1":          props["D1"],
                "D2":          props["D2"],
                "sigma_a1":    props["sigma_a1"],
                # ... etc
            }
\end{lstlisting}

%==============================================================================
\section{Lesson 4: 2D Neutron Diffusion Theory}
%==============================================================================

\subsection{From 1D to 2D}

\begin{physicsbox}{The Diffusion Equation in 2D}
In 1D, neutrons only move in x. In 2D, they move in both x and y.
Adding the y-direction leakage term:

\textbf{1D:}
\begin{equation}
-D \frac{d^2\phi}{dx^2} + \Sigma_a \phi = S
\end{equation}

\textbf{2D:}
\begin{equation}
-D\left(\frac{\partial^2\phi}{\partial x^2} +
        \frac{\partial^2\phi}{\partial y^2}\right) + \Sigma_a \phi = S
\end{equation}

The sum of second derivatives is the \textbf{Laplacian} -- it captures
total neutron leakage in all directions.
\end{physicsbox}

\subsection{Two-Group Equations}

\begin{physicsbox}{2-Group Diffusion System}
\textbf{Group 1 (Fast):}
\begin{equation}
-D_1\nabla^2\phi_1 + (\Sigma_{a1} + \Sigma_{s12})\phi_1
= \frac{1}{k}(\nu\Sigma_{f1}\phi_1 + \nu\Sigma_{f2}\phi_2)
\end{equation}

\textbf{Group 2 (Thermal):}
\begin{equation}
-D_2\nabla^2\phi_2 + \Sigma_{a2}\phi_2 = \Sigma_{s12}\phi_1
\end{equation}

Note: $\Sigma_{s12}$ appears in Group 1 as \textbf{removal} (loss) and
in Group 2 as \textbf{source} (gain).
\end{physicsbox}

\subsection{The 5-Point Stencil}

\begin{physicsbox}{Finite Difference Approximation}
At cell $(i,j)$, the Laplacian is approximated using 4 neighbors:
\begin{equation}
\nabla^2\phi \approx
\frac{\phi_{i+1,j} - 2\phi_{i,j} + \phi_{i-1,j}}{\Delta x^2}
+ \frac{\phi_{i,j+1} - 2\phi_{i,j} + \phi_{i,j-1}}{\Delta y^2}
\end{equation}

The center coefficient (all removal terms):
\begin{equation}
C_1 = \frac{2D_1}{\Delta x^2} + \frac{2D_1}{\Delta y^2} + \Sigma_{a1} + \Sigma_{s12}
\end{equation}

For the IAEA benchmark with $D_1=1.5$, $\Sigma_{a1}=0.01$,
$\Sigma_{s12}=0.02$, $\Delta x=1$ cm:
\begin{equation}
C_1 = 2(1.5) + 2(1.5) + 0.01 + 0.02 = 6.03 \text{ cm}^{-1}
\end{equation}
\end{physicsbox}

%==============================================================================
\section{Lesson 5: Matrix Assembly and Power Iteration}
%==============================================================================

\subsection{Matrix Structure}

\begin{physicsbox}{The Eigenvalue Problem}
For total unknowns $N = n_x \times n_y$ per group, with 2 groups:
\begin{equation}
\mathbf{A}\boldsymbol{\phi} = \frac{1}{k}\mathbf{F}\boldsymbol{\phi}
\end{equation}

Matrix $\mathbf{A}$ (size $2N \times 2N$): leakage + removal\\
Matrix $\mathbf{F}$ (size $2N \times 2N$): fission source\\
Vector $\boldsymbol{\phi}$ (size $2N$): flux [group1 cells, group2 cells]
\end{physicsbox}

\begin{lstlisting}[caption={Matrix Assembly -- The Core of the Solver}]
def _build_matrices(self):
    for j in range(ny):         # outer loop: y-direction (rows)
        for i in range(nx):     # inner loop: x-direction (cols)
            n  = j*nx + i       # 1D cell index

            xs = self._get_xs(i, j)  # cross sections at this cell
            D1, D2     = xs["D1"], xs["D2"]
            siga1      = xs["sigma_a1"]
            siga2      = xs["sigma_a2"]
            sigs12     = xs["sigma_s12"]
            nusigf1    = xs["nu_sigma_f1"]
            nusigf2    = xs["nu_sigma_f2"]

            Ex1 = D1/dx**2;  Ey1 = D1/dy**2
            Ex2 = D2/dx**2;  Ey2 = D2/dy**2

            C1 = siga1 + sigs12  # start with physics terms
            C2 = siga2

            # West neighbor
            if i > 0:
                add_A(n, j*nx+(i-1), -Ex1)  # coupling to left
                add_A(n+N, j*nx+(i-1)+N, -Ex2)
                C1 += Ex1;  C2 += Ex2
            else:
                C1 += self._boundary_coeff(D1, dx, "left")
                C2 += self._boundary_coeff(D2, dx, "left")

            # (East, South, North similarly...)

            add_A(n,   n,   C1)    # diagonal LAST (after all neighbors)
            add_A(n+N, n+N, C2)

            # Scattering: group 1 -> group 2
            add_A(n+N, n, -sigs12)

            # Fission source
            add_F(n,   n,   nusigf1)
            add_F(n,   n+N, nusigf2)
\end{lstlisting}

\begin{warnbox}{Critical Bug Pattern: Diagonal Added Too Early}
The center coefficient C1 accumulates leakage terms from each neighbor.
If \texttt{add\_A(n, n, C1)} is called BEFORE processing neighbors,
the diagonal will be wrong. Always add diagonal LAST.
\end{warnbox}

\subsection{Boundary Conditions}

\begin{lstlisting}[caption={Extrapolated Boundary Condition}]
def _boundary_coeff(self, D, h, side):
    bc_type = self.bc.get(side, "vacuum")

    if bc_type == "reflective":
        return 0.0                     # no leakage at symmetry plane

    elif bc_type == "vacuum":
        # Extrapolated BC: flux = 0 at d_ex = 2.1312*D beyond boundary
        d_ex = 2.1312 * D
        return D / (h * (h/2 + d_ex)) # less leakage than simple BC

    else:
        return D / h**2                # simple zero-flux BC
\end{lstlisting}

\begin{lessonbox}{Why Reflective BC for Quarter-Core?}
The IAEA benchmark is a quarter-core model. The symmetry planes (left and
bottom boundaries) represent the CENTER of the reactor.
At the center, $\partial\phi/\partial n = 0$ (zero current = no net flow).
Using \textbf{vacuum} BC at these planes would force $\phi=0$ at core center
-- completely wrong! This was the root cause of k-eff = 0.888 before the fix.
After adding reflective BC: k-eff = 0.998.
\end{lessonbox}

\subsection{Power Iteration}

\begin{lstlisting}[caption={Power Iteration -- Eigenvalue Solver}]
def solve(self):
    phi = np.ones(2 * N)   # initial flat flux guess
    k   = 1.0              # initial k-eff guess

    self._build_matrices()

    for iteration in range(max_iter):
        # Compute fission source
        fission_source = self.F.dot(phi)

        # Solve linear system: A * phi_new = (1/k) * F * phi
        rhs     = fission_source / k
        phi_new = linalg.spsolve(self.A, rhs)

        # Update k-eff (ratio of new to old fission source)
        fission_new = self.F.dot(phi_new)
        k_new = k * np.sum(fission_new) / np.sum(fission_source)

        # Convergence check
        k_error   = abs(k_new - k)
        phi_error = np.max(np.abs(phi_new - phi)/(np.abs(phi_new)+1e-10))

        phi = phi_new / np.max(phi_new)  # normalize
        k   = k_new

        if k_error < tol and phi_error < tol:
            break

    # Reshape 1D vector to 2D arrays
    self.flux1 = phi[:N].reshape(ny, nx)
    self.flux2 = phi[N:].reshape(ny, nx)
    self.k_eff = k
\end{lstlisting}

%==============================================================================
\section{Lesson 6: Multi-Group Solver (SolverMG)}
%==============================================================================

\subsection{Generalizing from 2-Group to G-Group}

\begin{physicsbox}{G-Group Center Coefficient}
For group $g$ ($g = 0, 1, \ldots, G-1$), the removal cross section includes
downscattering to ALL lower groups:
\begin{equation}
\Sigma_{r,g} = \Sigma_{a,g} + \sum_{g' > g} \Sigma_{s,g \to g'}
\end{equation}

Example for 4-group fuel ($g=1$):
\begin{equation}
C_2 = \frac{2D_2}{\Delta x^2} + \frac{2D_2}{\Delta y^2}
      + \Sigma_{a2} + \Sigma_{s,2\to3} + \Sigma_{s,2\to4}
\end{equation}
\end{physicsbox}

\begin{lstlisting}[caption={G-Group Cell Index and Matrix Assembly}]
def _cell_index(self, g, i, j):
    """Convert (group, i, j) to global matrix index."""
    n = j * self.nx + i      # spatial cell index
    return g * self.N + n    # group g starts at g*N

# For G=4, N=2500:
# Group 0: indices 0    to 2499
# Group 1: indices 2500 to 4999
# Group 2: indices 5000 to 7499
# Group 3: indices 7500 to 9999

def _build_matrices(self):
    for j in range(self.ny):
        for i in range(self.nx):
            xs = self._get_xs(i, j)
            for g in range(self.G):
                row = self._cell_index(g, i, j)

                # Center coefficient with all downscattering
                C = sigma_a_eff[g] + sum(
                    sigma_s[g][g2]
                    for g2 in range(self.G) if g2 != g
                )

                # Spatial neighbors (same group g, different i/j)
                if i > 0:
                    add_A(row, self._cell_index(g, i-1, j), -Dx)
                    C += Dx

                # Downscattering (same cell i,j, different group g2)
                for g_from in range(self.G):
                    if g_from == g: continue
                    add_S(row, self._cell_index(g_from,i,j),
                          sigma_s[g_from][g])

                # Fission source with chi spectrum
                for g2 in range(self.G):
                    add_F(row, self._cell_index(g2,i,j),
                          chi[g] * nusigf[g2])

                add_A(row, row, C)  # diagonal last!
\end{lstlisting}

\begin{keybox}{Key Insight: Spatial vs Group Coupling}
\begin{itemize}
  \item \textbf{Spatial neighbor}: same group $g$, different $(i,j)$
        -- changes position, keeps energy
  \item \textbf{Scattering}: different group $g'$, same $(i,j)$
        -- changes energy, keeps position
  \item \textbf{Fission}: different group $g'$, same $(i,j)$
        -- neutron born in group determined by chi
\end{itemize}
\end{keybox}

%==============================================================================
\section{Lesson 7: IAEA 2D PWR Benchmark}
%==============================================================================

\subsection{Benchmark Description}

\begin{lessonbox}{IAEA Benchmark Geometry}
Quarter-core model ($170 \times 170$ cm) with three fuel material types and
a water reflector. Each assembly is $10 \times 10$ cm. The staircase layout:

\medskip
\begin{center}
\begin{tabular}{ccccccccc}
F1&F1&F1&F1&F1&F1&F1&F1&R\\
F1&F1&F1&F1&F1&F1&F1&F2&R\\
F1&F1&F1&F1&F1&F2&F2&R&R\\
F1&F1&F1&F1&F2&F2&R&R&R\\
F1&F1&F1&F2&F2&CR&R&R&R\\
F1&F1&F2&F2&CR&R&R&R&R\\
F1&F1&F2&R&R&R&R&R&R\\
F1&F2&R&R&R&R&R&R&R\\
R&R&R&R&R&R&R&R&R\\
\end{tabular}
\end{center}

\medskip
Boundary: \textbf{reflective} on left and bottom (symmetry planes),
\textbf{vacuum} on right and top.
\end{lessonbox}

\subsection{Cross Sections Used}

\begin{center}
\begin{tabular}{lrrrrrr}
\toprule
Material & $D_1$ & $D_2$ & $\Sigma_{a1}$ & $\Sigma_{a2}$
         & $\Sigma_{s12}$ & $\nu\Sigma_{f2}$ \\
\midrule
Fuel 1   & 1.500 & 0.400 & 0.0100 & 0.0850 & 0.0200 & 0.135 \\
Fuel 2   & 1.500 & 0.400 & 0.0100 & 0.1300 & 0.0200 & 0.135 \\
Fuel+CR  & 1.500 & 0.400 & 0.0100 & 0.1800 & 0.0200 & 0.135 \\
Reflector& 2.000 & 0.300 & 0.0000 & 0.0100 & 0.0400 & 0.000 \\
\bottomrule
\end{tabular}
\end{center}

\subsection{Results and Error Analysis}

\begin{center}
\begin{tabular}{lrr}
\toprule
Parameter & Value & Notes \\
\midrule
ANGKOR k-eff & 0.99837 & \\
Reference k-eff & 1.02959 & IAEA published \\
Error & -3086 pcm & Expected \\
Iterations & 73 & Power iteration \\
Runtime & $\sim$43 seconds & Python/SciPy \\
Mesh convergence & 1.1 pcm & 170 vs 340 mesh \\
\bottomrule
\end{tabular}
\end{center}

\begin{physicsbox}{Error Decomposition}
The -3086 pcm error is physically expected:
\begin{itemize}
  \item $\sim$2000 pcm: Diffusion approximation vs transport theory (P1 vs B/N)
  \item $\sim$1000 pcm: 2-group energy structure (spectral effects at fuel-reflector)
  \item $<$100 pcm: Numerical (mesh, boundary) -- verified by convergence study
\end{itemize}
\textbf{Conclusion}: No code bugs. This is the correct answer for this level of approximation.
\end{physicsbox}

%==============================================================================
\section{Lesson 8: Cross Section Generation with OpenMC}
%==============================================================================

\subsection{Group Condensation Theory}

\begin{physicsbox}{Flux-Weighted Average}
Cross sections must be averaged over energy groups using flux as a weight,
because reactions occur where neutrons actually are:
\begin{equation}
\bar{\sigma}_g = \frac{\displaystyle\int_{E_{g+1}}^{E_g}\sigma(E)\phi(E)\,dE}
                      {\displaystyle\int_{E_{g+1}}^{E_g}\phi(E)\,dE}
\end{equation}
Simple (unweighted) averages give wrong results because they ignore that
thermal neutrons dominate reaction rates.
\end{physicsbox}

\begin{lstlisting}[caption={Group Condensation Implementation}]
def condense(self, nuclide_data, group_structure="2-group"):
    boundaries = self.Groups_Structure[group_structure]
    G          = len(boundaries) - 1

    E     = nuclide_data["energy"]
    sig_f = nuclide_data["fission"]
    phi   = 1.0 / E          # 1/E slowing-down spectrum

    results = []
    for g in range(G):
        E_max = boundaries[g]      # high energy boundary
        E_min = boundaries[g+1]    # low energy boundary

        mask    = (E >= E_min) & (E <= E_max)
        E_g     = E[mask]
        sig_f_g = sig_f[mask]
        phi_g   = phi[mask]

        numer = np.trapezoid(sig_f_g * phi_g, E_g)  # numerator
        denom = np.trapezoid(phi_g, E_g)             # denominator
        results.append(numer / denom if denom > 0 else 0.0)

    return results
\end{lstlisting}

\begin{keybox}{Numerical Convergence Study}
\begin{center}
\begin{tabular}{rrr}
\toprule
Grid Points & G2 $\sigma_f$ (thermal) & Change \\
\midrule
100   & 5397.4 barns & -- \\
1000  & 5266.4 barns & 2.4\% \\
10000 & 5262.7 barns & 0.06\% \\
\bottomrule
\end{tabular}
\end{center}
Lesson: 100 points is NOT sufficient for nuclear data integration.
10,000+ points required for $<$0.1\% accuracy.
\end{keybox}

\subsection{OpenMC Pin Cell Workflow}

\begin{lstlisting}[caption={OpenMC PWR Pin Cell Setup}]
import openmc

# Materials
uo2 = openmc.Material(name="UO2 Fuel")
uo2.set_density("g/cm3", 10.29)
uo2.add_nuclide("U235", 3.1,   percent_type="ao")  # 3.1% enriched
uo2.add_nuclide("U238", 96.9,  percent_type="ao")
uo2.add_nuclide("O16",  200.0, percent_type="ao")  # 2 O per U atom

water = openmc.Material(name="Water")
water.set_density("g/cm3", 0.7)
water.add_nuclide("H1",  2.0, percent_type="ao")
water.add_nuclide("O16", 1.0, percent_type="ao")
water.add_s_alpha_beta("c_H_in_H2O")  # thermal scattering law

# Geometry: fuel pin + water
fuel_radius = openmc.ZCylinder(r=0.4096)  # cm
pitch = 1.26  # cm

fuel_cell  = openmc.Cell(fill=uo2,   region=-fuel_radius)
water_cell = openmc.Cell(fill=water, region=+fuel_radius)

# Settings: 150 batches, 50 inactive, 10000 particles
settings = openmc.Settings()
settings.batches   = 150
settings.inactive  = 50    # let fission source converge first
settings.particles = 10000
\end{lstlisting}

\begin{lessonbox}{Result: k-inf = 1.394 for Fresh Fuel Pin}
This is physically correct: infinite lattice (no leakage) + fresh fuel
(no burnup, no fission products) + no control rods = high k-eff.
Real core k-eff $\approx$ 1.03 after leakage + control rods + feedback.
\end{lessonbox}

\begin{warnbox}{Lesson Learned: Two-Step Method Challenge}
Pin cell XS from infinite lattice $\neq$ assembly XS for finite core.
The core requires a \textbf{critical spectrum correction} (B$^2$ search)
applied during XS generation. Without it, 4-group k-eff diverges to 1.37+.
Proper XS generation needs a full assembly calculation in CASMO or Serpent.
\end{warnbox}

%==============================================================================
\section{Lesson 9: CMFD Acceleration (In Progress)}
%==============================================================================

\subsection{Motivation: Dominance Ratio}

\begin{physicsbox}{Why Power Iteration is Slow}
The number of iterations needed is governed by the dominance ratio:
\begin{equation}
DR = \frac{k_1}{k_0} \quad \text{(ratio of 1st to fundamental eigenvalue)}
\end{equation}
For our reactor: $DR \approx 0.99$ means k-eff and the first harmonic are
nearly identical. Power iteration needs many iterations to separate them.
CMFD accelerates by solving a coarse mesh problem that has a smaller
dominance ratio.
\end{physicsbox}

\subsection{CMFD Algorithm}

\begin{lstlisting}[caption={CMFD Flux Homogenization}]
class CMFD:
    def __init__(self, solver, rf=10):
        if solver.nx % rf != 0:
            raise ValueError(f"nx={solver.nx} not divisible by rf={rf}")
        if solver.ny % rf != 0:
            raise ValueError(f"ny={solver.ny} not divisible by rf={rf}")

        self.rf   = rf
        self.nx_c = solver.nx // rf   # coarse cells in x
        self.ny_c = solver.ny // rf   # coarse cells in y
        self.Dx   = solver.dx * rf    # coarse cell width
        self.Dy   = solver.dy * rf    # coarse cell height

    def homogenize_flux(self, phi1_fine, phi2_fine):
        """Average fine mesh flux onto coarse mesh."""
        phi1_c = np.zeros((self.ny_c, self.nx_c))
        phi2_c = np.zeros((self.ny_c, self.nx_c))

        for J in range(self.ny_c):
            for I in range(self.nx_c):
                j_start = J * self.rf
                j_end   = (J+1) * self.rf   # no -1 (Python exclusive)
                i_start = I * self.rf
                i_end   = (I+1) * self.rf

                # numpy slicing: [rows, cols] = [j, i]
                phi1_c[J][I] = np.mean(phi1_fine[j_start:j_end,
                                                   i_start:i_end])
                phi2_c[J][I] = np.mean(phi2_fine[j_start:j_end,
                                                   i_start:i_end])
        return phi1_c, phi2_c
\end{lstlisting}

\begin{physicsbox}{d-hat Correction Coefficient}
At each coarse mesh interface, the correction coefficient ensures that
the coarse mesh current matches the fine mesh current:
\begin{equation}
\hat{d} = \frac{J_{\text{fine}} -
               \left(-\tilde{D}\dfrac{\Delta\Phi}{\Delta x}\right)}{\bar{\Phi}}
\end{equation}
The interface diffusion coefficient uses the \textbf{harmonic mean}
(not simple average) to correctly handle material discontinuities:
\begin{equation}
D_{\text{interface}} = \frac{2D_i D_{i+1}}{D_i + D_{i+1}}
\end{equation}
Simple average would overestimate current at fuel-reflector boundaries.
\end{physicsbox}

%==============================================================================
\section{Lesson 10: Software Engineering Practices}
%==============================================================================

\subsection{Key Python Concepts Mastered}

\begin{lstlisting}[caption={Important Python Patterns Used Throughout}]
# 1. context manager (always use for files)
with open("file.yaml", "r", encoding="utf-8") as f:
    data = yaml.safe_load(f)

# 2. safe dictionary lookup with default
sol.get("max_iterations", 1000)  # returns 1000 if key missing

# 3. list comprehension (cell centers)
x_centers = np.array([(i + 0.5) * dx for i in range(nx)])

# 4. boolean mask (numpy)
mask = (E >= E_min) & (E <= E_max)  # True where condition met
E_g  = E[mask]                       # select matching elements

# 5. sparse matrix construction
A = sparse.csr_matrix((vals, (rows, cols)), shape=(size, size))

# 6. f-string formatting
print(f"k-eff = {k:.6f}  error = {dk:.2e}")

# 7. try/except for validation testing
try:
    bad = Region("bad", 80.0, 20.0, 0.0, 100.0, "water")
except ValueError as e:
    print(f"Caught correctly: {e}")
\end{lstlisting}

\subsection{Common Bug Patterns}

\begin{center}
\begin{longtable}{p{3.5cm} p{5cm} p{5cm}}
\toprule
\textbf{Bug Type} & \textbf{Wrong} & \textbf{Correct} \\
\midrule
Loop bounds swapped & \texttt{for j in range(nx)} & \texttt{for j in range(ny)} \\
Wrong property name & \texttt{region.material\_map} & \texttt{region.material} \\
typo & \texttt{np.meaen()} & \texttt{np.mean()} \\
Wrong axis & \texttt{phi[i\_start:i\_end, j...]} & \texttt{phi[j\_start:j\_end, i...]} \\
Diagonal too early & \texttt{add\_A(n,n,C); C+=Ex} & \texttt{C+=Ex; add\_A(n,n,C)} \\
Hardcoded instead of variable & \texttt{self.Groups["2-group"]} & \texttt{self.Groups[group\_structure]} \\
Off-by-one in slice & \texttt{i\_end = (I+1)*rf - 1} & \texttt{i\_end = (I+1)*rf} \\
\bottomrule
\end{longtable}
\end{center}

\subsection{Git Workflow}

\begin{lstlisting}[language=bash, caption={Standard Git Workflow}]
# Golden rule: ALWAYS pull before push
git pull origin main
git add <specific files>   # never use "git add ." carelessly
git commit -m "feat: description of what was built"
git push

# Check what is tracked:
git ls-files

# Remove file from tracking (keep local copy):
git rm --cached filename

# See commit history:
git log --oneline
\end{lstlisting}

\begin{lessonbox}{Repository Philosophy}
Only commit what you can defend:
\begin{itemize}
  \item Validated results only (not work-in-progress with known wrong answers)
  \item Learning files go in .gitignore
  \item OpenMC experiment files (.xml, statepoint.h5) go in .gitignore
  \item Unvalidated cross section files go in .gitignore
\end{itemize}
\end{lessonbox}

%==============================================================================
\section{Roadmap for Future Sessions}
%==============================================================================

\subsection{Next Immediate Steps}

\begin{enumerate}
  \item \textbf{Complete CMFD:} Implement \texttt{compute\_dhat()},
        \texttt{build\_coarse\_matrices()}, \texttt{solve\_coarse()},
        \texttt{rebalance\_flux()}, and \texttt{accelerate()}.
        Expected: 5$\times$ speedup on IAEA benchmark.

  \item \textbf{Test CMFD:} Confirm same k-eff with and without acceleration.
        Measure iteration count and runtime reduction.

  \item \textbf{Proper 4-group XS:} Run full assembly calculation in OpenMC
        with B$^2$ search to get transport-corrected few-group XS.
        This requires a $17\times17$ pin cell lattice calculation, not just
        a single pin.

  \item \textbf{Theory Manual:} Document all physics and numerical methods
        in \texttt{docs/ANGKOR\_Theory\_Manual.tex}.
\end{enumerate}

\subsection{Medium-Term Goals}

\begin{enumerate}
  \item Cylindrical geometry for fuel pins
  \item 3D Cartesian geometry extension
  \item Adjoint flux solver (reactivity worth)
  \item Burnup and depletion module
  \item OECD/NEA C5G7 benchmark (7-group)
\end{enumerate}

\subsection{Long-Term Vision}

\begin{lessonbox}{Cambodia's Nuclear Engineering Legacy}
ANGKOR's goal is not just to compute k-eff. It is to demonstrate that
Cambodia can contribute to the global nuclear engineering community
through open-source software -- without relying on expensive commercial
codes. The code should be readable, modifiable, and useful to any
graduate student or researcher worldwide.

The ITC group is also exploring ML surrogate models for SMR k-effective
prediction. ANGKOR provides the reference data generation capability
for such work.
\end{lessonbox}

%==============================================================================
\vfill
\begin{center}
\rule{0.5\linewidth}{0.5pt}\\[0.5cm]
{\small ANGKOR -- Advanced Neutron Group-diffusion K-eigenvalue Of Reactors}\\
{\small\textit{Developed at the Institute of Technology of Cambodia}}\\
{\small Phnom Penh, Cambodia}\\
{\small \url{https://github.com/muthboravy007/ANGKOR}}\\[0.3cm]
{\small\itshape ``Named after Angkor Wat -- because Cambodia's scientific legacy continues.''}
\end{center}

\end{document}
